\chapter{Concluzii}

Lucrarea a propus și a evaluat un pipeline complet, executabil local, pentru un asistent conversațional dedicat procedurilor administrative românești, ancorat în datele publicate de Code for Romania. Față de lucrările grupului de cercetare (Gheorghe et al., 2026 \cite{gheorghe2026robust}; Paduraru et al., 2026 \cite{paduraru2026guardrail}), contribuțiile principale sunt: o sursă de date și evenimente de viață disjuncte; un strat de verificare a ieșirii (șase reguli evaluate, plus R7 pentru supra-refuz) care implementează verificarea răspunsului identificată în literatură ca lucrare viitoare; și o arhitectură integral locală, pe hardware de larg consum.

Evaluarea cantitativă susține valoarea contribuțiilor: pe 3B, \textbf{67\%} dintre răspunsuri ar fi fost servite cu defect fără stratul de verificare (în majoritate din lipsa citațiilor), iar reîncercarea le repară sau le reține controlat. Trecerea la modelul mai mare reduce documentele inventate de la 26 de cazuri la 2 (o reducere de peste 90\%), validând un mecanism de verificare independent de model.

Un rezultat al evaluării clarifică o alegere de proiectare: studiul comparativ de cuantizare arată că principalul factor de degradare este precizia, nu scara. Configurația 3B-Q8 (74\% refusal accuracy) depășește 3B-Q4 (62\%) și întrece modele mai mari, dar slab cuantizate (7B-Q4 69\%, 14B-Q4 66\%), deci un model mic, bine cuantizat, este competitiv pentru rularea locală. Analiza modurilor de eșec a condus și la regula R7 de detecție a supra-refuzului, care, calibrată pe setul de aur, detectează 73\% dintre refuzurile nejustificate fără fals pozitive.

Experimentul cu un model specializat pe limba română (RoLlama3-8B) confirmă o direcție și deschide alta: specializarea reduce drastic refuzul nejustificat (3 cazuri, față de 33), dar deplasează eroarea către sub-refuz și o citare mai slabă, ceea ce justifică valoarea stratului de verificare. Două direcții rămân prioritare pentru dezvoltarea ulterioară: o regăsire conștientă de intenție, pentru întrebările care combină mai multe evenimente, și o regulă suplimentară R8 de ancorare cross-eveniment, integrabilă în stratul de verificare fără modificarea logicii de generare.

\section{Discuție și limitări}
\label{sec:limitari}

Verificarea adresează eficient eșecul \emph{zgomotos} (inventarea de documente), însă modul de eșec \emph{dominant}, refuzul nejustificat, rămâne în mare parte invizibil pentru R1-R6 și este doar semnalat de R7. Prin urmare, valoarea verificării este reală, dar parțială.

Câteva limitări țin de \emph{validitatea de construct}: breadcrumb recall@k și lexical keyword recall sunt \emph{proxy-uri}, triangulate cu o evaluare manuală (15 cazuri), cu un judecător LLM și cu auditul R4 ($\approx$ 90\% precizie), iar cuvintele-cheie au fost validate automat ca prezente în corpusul-sursă; în plus, ,,eroarea nesemnalată" este o definiție operațională (declanșarea unei reguli), nu o judecată independentă.

Din perspectiva \emph{validității interne}, recall-ul constant izolează diferențele la nivelul generării, însă pragul R7 este calibrat și evaluat pe același set de 98 de cazuri, deci performanța de detecție este \emph{in-sample}, validarea pe un set ținut deoparte rămânând lucrare viitoare; întrucât toate configurațiile au fost evaluate pe aceeași infrastructură, diferențele de calitate sunt atribuibile exclusiv modelului.

În ceea ce privește \emph{validitatea externă}, setul de aur (98 de cazuri, generat asistat de LLM) acoperă doar trei evenimente de viață, un singur model de embedding și exclusiv limba română, pe un instantaneu al corpusului, astfel încât generalizarea la alte domenii sau limbi nu este garantată.
